\chapter{Preliminary Research}
\label{ch:research}

\section{Literature Review}

\subsection{Presentation attacks, deepfakes, and injection}

PAD traditionally evaluates whether a biometric sample is bona fide or a
presentation attack at the capture device. ISO/IEC 30107-3 defines principles for
testing and reporting PAD performance, including attack-type classification, but
explicitly excludes attacks outside the capture device from its scope
\cite{iso30107}. This boundary matters: a detector that recognises screen moir\'e
or print texture does not establish that a video stream came from a genuine
camera. ENISA's later guidance therefore treats injection attack detection (IAD)
alongside PAD and recommends a multi-layer client--server approach
\cite{enisa2024}. NIST SP 800-63A-4 similarly calls for confidence in the genuine
sensor, protected communications, manipulated-media analysis, and documented
manual processes \cite{nist80063a4}.

FaceForensics++ provides a controlled benchmark with several facial manipulation
methods and compression levels, and established Xception as a strong baseline
\cite{rossler2019faceforensics,chollet2017xception}. Celeb-DF-v2 was designed with
higher-quality face swaps and fewer conspicuous synthesis artefacts, making it a
more challenging test of shortcut-dependent detectors \cite{li2020celebdf}.
Cross-dataset deepfake work repeatedly shows that strong in-domain performance
does not guarantee unseen-generator performance; video-oriented approaches such
as AltFreezing were developed specifically to improve this transfer
\cite{wang2023altfreezing}.

\subsection{Temporal and domain-generalised face anti-spoofing}

MobileNetV3 uses hardware-aware architecture search and efficient blocks intended
for resource-constrained inference \cite{howard2019mobilenetv3}. In this project,
its per-frame features are averaged over a short clip, providing a simple temporal
baseline without the cost of full spatiotemporal attention.

$G^2V^2$former combines cropped face appearance with facial landmarks. It
factorises spatial and temporal attention, introduces Kronecker temporal
attention, and uses landmark motion to guide facial-expression features
\cite{yang2024g2v2}. It is therefore relevant when motion contains evidence that
is weak or ambiguous in any individual frame.

GD-FAS addresses a different failure mode. Its group-wise scaling risk
minimisation and feature orthogonal decomposition aim to align both the direction
and bias of local decision boundaries while separating domain-invariant from
domain-specific features \cite{jung2025gdfas}. It remains a frame-level spatial
model in this project; comparing it with $G^2V^2$former helps separate the value
of temporal structure from the value of domain-generalised representation.

\begin{table}[H]
\centering
\caption{Research approaches and their roles in the project.}
\label{tab:literature-map}
\begin{tabularx}{\textwidth}{@{}p{28mm}p{27mm}Y Y@{}}
\toprule
Approach & Primary evidence & Strength & Limitation relevant here \\
\midrule
MobileNetV3 pooling & Short face clip & Efficient mobile temporal summary & Mean
pooling can discard ordering and subtle motion. \\
$G^2V^2$former & Face video and 68 landmarks & Explicit photometric--dynamic
fusion & Landmark extraction and attention increase latency and failure modes. \\
GD-FAS & Single face frames and domain labels & Explicit domain-invariant feature
objective & Does not directly model temporal consistency. \\
Universal Xception & Single face frame & Simple common detector and baseline & A
single boundary may average incompatible manipulation artefacts. \\
Routed specialists & Embedding plus expert bank & Allocates capacity by
manipulation pattern & Router errors, training leakage, and compute multiply risk. \\
Capture/IAD controls & Device and channel signals & Addresses camera bypass & Not
implemented by a media classifier alone. \\
\bottomrule
\end{tabularx}
\end{table}

\section{Existing Solutions and Technologies}

Commercial liveness systems commonly advertise passive analysis, active
challenge--response, or combinations of both. Active systems ask the user to
perform a random action; they may provide session-specific evidence but add
friction and accessibility concerns. Passive systems analyse ordinary capture in
the background; they can improve usability but must not be treated as proof of a
trusted source. Vendor claims are useful indicators of product direction, not
independent evidence of accuracy.

The technological alternatives considered were lightweight convolutional
networks, video transformers, domain-generalisation methods, universal deepfake
detectors, specialist ensembles, active challenges, and capture-channel controls.
No one option covers every threat. The selected design therefore treats the
classifier as one layer in a risk-based verification process.

\section{Market and User Research}

Pathao supplied the initial financial-verification problem and early context, but
the team did not receive private operational data, a grant, or sustained product
collaboration. Consequently, this report does not claim direct Pathao user
research, production integration, or measured business impact. The requirements
instead derive from the initial problem framing, security guidance, public
literature, and prototype observation.

The affected user groups extend beyond fraud analysts. Legitimate applicants need
a quick, comprehensible process that works across devices, connectivity levels,
skin tones, ages, abilities, and religious or cultural presentation. Reviewers
need bounded workloads and enough context to avoid blind trust in scores. A
financial provider needs auditable decisions, recoverable failure handling, and
controls against both fraud and unlawful or excessive biometric processing.

Before deployment, the market/user study must therefore add interviews or tests
with applicants, accessibility participants, fraud reviewers, support personnel,
security engineers, privacy/legal personnel, and product owners. Required outputs
include acceptable capture time, retry limits, low-bandwidth behaviour, consent
language, review turnaround, and an explicit false-accept/false-reject cost model.

\section{Feasibility Study}

\textbf{Technical feasibility.} A temporal MobileNetV3 has been exported and
integrated into a Flutter capture path, showing that local inference is feasible.
The stronger models run in research/server environments. End-to-end selective
escalation and capture attestation remain unimplemented, so system feasibility is
only partially demonstrated.

\textbf{Data feasibility.} Public datasets make repeatable research possible, but
they cannot represent the final operating population or attack distribution.
Licences, consent conditions, identity leakage, demographic composition, and
subject-disjoint splits must be audited before reuse. No private Pathao data is
available for deployment-domain validation.

\textbf{Operational feasibility.} A cascade can reduce server load by retaining a
fast local path, but only after scores are calibrated and thresholds are selected
from operational costs. Network failure, poor capture, device compromise, and
review capacity require explicit fallback behaviour.

\textbf{Economic and legal feasibility.} Open-source tooling and public data
reduce prototype cost, whereas GPU training, secure infrastructure, testing,
monitoring, and human review remain material. NIST and ENISA are authoritative
engineering guidance, not declarations of legal compliance in Bangladesh. Local
privacy, financial, cyber-security, consumer-protection, and evidentiary
requirements require review by qualified stakeholders before deployment.

\textbf{Conclusion.} Research prototyping is feasible, and the on-device stage
has been demonstrated. Production adoption is conditional on integrated security
controls, field data, independent PAD/IAD testing, subgroup analysis, operational
calibration, and governance.

